\documentclass[a4paper, 11pt]{article}
\usepackage{geometry}
\usepackage{graphicx}
\usepackage{amsmath}
\usepackage{float}
\usepackage{booktabs}

\geometry{margin=1in}

\begin{document}
	
	\title{Deep Learning Enhanced Virtual Sensing for Water Quality Estimation}
	\maketitle
	
	Efficient monitoring of water quality parameters is crucial for the successful management of aquaculture systems, ensuring the health and growth of aquatic life. Traditional sensor-based systems are often cost-prohibitive, limiting their adoption, particularly in resource-constrained settings. In this research, we introduce a cost-effective method that leverages virtual sensing to estimate dissolved oxygen (DO) levels using solubility and interdependency equations. Furthermore, we propose deep learning-enhanced virtual sensing (DLEVS) by employing deep recurrent neural networks (DL-RNN) to accurately estimate DO. 
	
	\section*{Virtual Sensing Using Solubility Equations}
	
	Our virtual sensing (VS) methodology leverages established solubility equations to estimate water quality parameters, with a specific focus on dissolved oxygen (DO). By utilizing mathematical models derived from comprehensive datasets, our VS system offers a cost-effective and practical alternative to physical sensors. The flexibility of this approach allows us to estimate DO concentrations irrespective of the source of water, making it applicable to various water bodies.
	
	\subsection*{VS Capabilities Irrespective of Water Source}
	
	The VS system, based on solubility equations, is designed to estimate water quality parameters across diverse water sources. Through mathematical models, the system predicts parameters that are challenging to measure directly through physical sensors. This flexibility ensures accurate estimations for a wide range of scenarios.
	
	\section*{Deep Learning Enhancement for Accuracy}
	
	To address the Controller's question regarding the learning mechanism and real-time error correction, our VS system incorporates deep learning recurrent neural network (DL-RNN) models through our Dissolved Oxygen Enhanced Virtual Sensing (DLEVS) mechanism. This mechanism aims to compensate for potential inaccuracies in the VS-based DO estimation, particularly when dealing with new, untrained water bodies in real time.
	
	\subsection*{DL-RNN Training and Error Prediction}
	
	The DL-RNN model undergoes comprehensive training on various water quality datasets to predict errors in VS-based DO estimates. This rigorous training process ensures the adaptability and high predictive accuracy of the model across diverse water sources. The utilization of multiple datasets contributes to the robustness of the DL-RNN, allowing it to capture the intricacies and variations present in different water quality scenarios.
	
%	\subsection*{Generalization of DL-RNN for Error Prediction}
%	
%	Importantly, the DL-RNN model is engineered to predict errors in virtual sensing without requiring an exhaustive study of the specific dataset characteristics. This unique attribute enables the model to generalize its learning, making it exceptionally versatile across a spectrum of water quality scenarios. By focusing on error prediction, rather than being dependent on a specific dataset, the DL-RNN ensures adaptability to new water bodies in real time, enhancing the system's overall effectiveness.
%	
%	\subsection*{Real-Time Measurements and DL Enhancement}
%	
%	In real-time scenarios, salinity and temperature are measured, and the VS system estimates DO using solubility equations. The DL-RNN model predicts errors based on the dynamic characteristics of the specific water source. By subtracting the predicted error from the VS-estimated DO, we obtain an enhanced DO estimate that is more accurate and reflective of real-time conditions.
	
	\subsection*{Real-Time Measurements and DL Enhancement}
	
	In real-time scenarios, salinity and temperature are measured, and the VS system estimates DO using solubility equations. The DL-RNN model is engineered to predict errors in virtual sensing without requiring an exhaustive study of specific dataset characteristics. This unique attribute enables the model to generalize its learning, making it exceptionally versatile across a spectrum of water quality scenarios. By focusing on error prediction, rather than being dependent on a specific dataset, the DL-RNN ensures adaptability to new water bodies in real time, enhancing the system's overall effectiveness. The DL-RNN model predicts errors based on the dynamic characteristics of the specific water source. By subtracting the predicted error from the VS-estimated DO, we obtain an enhanced DO estimate that is more accurate and reflective of real-time conditions.
	
	
	\section*{Conclusion}
	
	In conclusion, our VS system, leveraging solubility equations and DL-RNN models, demonstrates the capability to estimate water quality parameters irrespective of the source of water. The integration of DL-RNN models through the DLEVS mechanism enhances the accuracy of DO estimation by predicting and compensating for potential errors. Real-time measurements of salinity and temperature, coupled with virtual sensing and DL enhancement, provide a comprehensive and accurate approach to water quality estimation.
	
\end{document}
